\documentclass[conference]{IEEEtran}
\usepackage[T1]{fontenc}
\usepackage{booktabs}
\usepackage{url}

\begin{document}

\title{How Do People Actually Prompt? Prompt Smells in 1,000 Real ChatGPT Requests}

\author{\IEEEauthorblockN{Author Name}
\IEEEauthorblockA{Institute of Information Technology\\
Email: author@example.org}}

\maketitle

\begin{abstract}
%ABSTRACT%
\end{abstract}

\begin{IEEEkeywords}
prompt engineering, prompt smells, large language models, WildChat, empirical study
\end{IEEEkeywords}

\section{Introduction}
Most advice on writing prompts comes from guides and from lab studies with set tasks.
We know much less about how people write prompts when nobody is watching.
WildChat~\cite{wildchat} closes part of that gap. It holds about a million real
conversations that users had with ChatGPT in exchange for free access.

In code, a \emph{smell} is a sign that something may be wrong, even if the program
still runs~\cite{fowler}. We borrow the idea for prompts. A prompt smell is a
pattern that makes a request harder for a model to answer well, such as a missing
goal, an unclear ``it'', or two rules that cannot both hold.

We ask two questions.
\textbf{RQ1:} How common are prompt smells in real English prompts, and which ones
come up most?
\textbf{RQ2:} Do smell rates change with prompt length and with the model the user
was talking to?

%CONTRIB%

\section{Related Work}
White et al.~\cite{white} collect prompt patterns that work, and the Prompt
Report~\cite{promptreport} surveys prompting methods. Both describe what good
prompts look like rather than what users really write. Zamfirescu-Pereira
et al.~\cite{johnny} watched non-experts write prompts and found they often
give up after one failed try and lean on human-style instructions. Our study looks
at the same problem at a larger scale, using logs instead of a lab. We use an
LLM as the labeler, following the LLM-as-a-judge line of work~\cite{judge}.

\section{Method}
\subsection{Data and sampling}
We used one of the six training shards of WildChat (\texttt{train-00003-of-00006},
88,238 conversations). To avoid taking only the earliest conversations in the
file, we kept a conversation when a hash of its id fell in the lowest 1\%. This
gives a fixed, repeatable sample spread across the whole shard. From it we took
every user turn tagged as English that was not empty, which gave 1,003 prompts,
and sent all of them to the labeler.

\subsection{Smell catalog}
We used ten smells, shown in Table~\ref{tab:freq}. They cover missing context
(vague prompts, unclear references, no stated output format), structure (one
prompt asking for many things, bloat, repetition), contradictions (conflicting
constraints), and tone (excess persona, loaded framing, wrong register).

\subsection{LLM labeling}
Each prompt was sent on its own to \texttt{gpt-oss-120b}~\cite{gptoss} through
the Groq API at temperature 0. The instructions list the ten smells with a short
definition, mark the prompt with fixed start and end markers, and say whether the
prompt is the first turn or a follow-up, so that a short follow-up such as
``make it shorter'' is not marked vague just because it leans on earlier turns.
The model had to return JSON with a smell name and a one-sentence reason for each
smell found. We checked every answer against a schema and asked again, up to three
attempts in all, when it did not fit. Close spellings of a smell name were mapped
to the catalog name. Prompts longer than 12,000 characters were cut, and the model
was told so. Prompts that still failed were logged and left out.

\section{Results}
%RESULTS%

\section{Discussion}
%DISCUSSION%

\section{Threats to Validity}
\textbf{Labels.} The labels come from one LLM and were not checked by people.
The model may be too strict or too lenient, and we did not measure agreement with
human raters. The counts should be read as estimates.
\textbf{Sample.} We used one shard, a 1\% sample and English only. WildChat
language tags are sometimes wrong; we saw a few Chinese prompts tagged as English.
WildChat users chose to share their chats for free access, so they may not
represent all ChatGPT users.
\textbf{Context.} Each prompt was judged alone, with only a first-turn or follow-up
flag. A prompt that looks vague alone may be clear in its chat.

\section{Conclusion}
%CONCLUSION%

\section*{Data and Code Availability}
The analysis tool, the sampling settings and the labeled output
(\texttt{wildchat\_1k\_smells.json}) are available from the authors.

\begin{thebibliography}{9}
\bibitem{wildchat} W.~Zhao, X.~Ren, J.~Hessel, C.~Cardie, Y.~Choi, and Y.~Deng,
``WildChat: 1M ChatGPT interaction logs in the wild,'' in \emph{Proc. ICLR}, 2024.
\bibitem{fowler} M.~Fowler, \emph{Refactoring: Improving the Design of Existing
Code}. Addison-Wesley, 1999.
\bibitem{white} J.~White \emph{et al.}, ``A prompt pattern catalog to enhance
prompt engineering with ChatGPT,'' arXiv:2302.11382, 2023.
\bibitem{promptreport} S.~Schulhoff \emph{et al.}, ``The prompt report: A
systematic survey of prompting techniques,'' arXiv:2406.06608, 2024.
\bibitem{johnny} J.~D. Zamfirescu-Pereira, R.~Y. Wong, B.~Hartmann, and Q.~Yang,
``Why Johnny can't prompt: How non-AI experts try (and fail) to design LLM
prompts,'' in \emph{Proc. CHI}, 2023.
\bibitem{judge} L.~Zheng \emph{et al.}, ``Judging LLM-as-a-judge with MT-Bench and
Chatbot Arena,'' in \emph{Proc. NeurIPS Datasets and Benchmarks}, 2023.
\bibitem{gptoss} OpenAI, ``gpt-oss-120b \& gpt-oss-20b model card,''
arXiv:2508.10925, 2025.
\end{thebibliography}

\end{document}
